\subsection{对称形式幂级数}

设 I 为一个集合， \(\mathbf{X} = (\mathbf{X}_i)_{i\in I}\) 为一族未定元， \(\mathbf{A}[[\mathbf{X}]]\) 为 \(\mathbf{X}_i\) 中的形式幂级数代数。对每个置换 \(\sigma \in \mathfrak{S}_I\) ，存在唯一的连续自同构 \(\varphi_{\sigma}\) （作用于代数 \(\mathbf{A}[[\mathbf{X}]]\) ），它把每个 \(\mathbf{X}_i\) 映为 \(\mathbf{X}_{\sigma (i)}\) （ \(i\in I\) ）（IV，第28页，命题4）；显然 \(\sigma \mapsto \varphi_{\sigma}\) 是从 \(\mathfrak{S}_I\) 到代数 \(\mathbf{A}[[\mathbf{X}]]\) 的连续自同构群中的同态。设 \(f\in \mathbf{A}[[\mathbf{X}]]\) 为一个形式幂级数；我们令 \(\sigma f = \varphi_{\sigma}(f)\) ，并称形式幂级数 \(f\) 是对称的，如果 \(\sigma f = f\) （对每个 \(\sigma \in \mathfrak{S}_I\) ）。对称形式幂级数构成 \(\mathbf{A}[[\mathbf{X}]]\) 的一个闭子代数，记为 \(\mathbf{A}[[\mathbf{X}]]^{\mathrm{sym}}\) ，并配备由 \(\mathbf{A}[[\mathbf{X}]]\)

诱导的拓扑。设 T 为一个未定元。在形式幂级数环 A{[}{[}X, T{]}{]} 中，族 \((\mathbf{X}_i\mathbf{T})_{i\in I}\) 是可求和的，因而族 \((1 + \mathbf{X}_i\mathbf{T})_{i\in I}\) 是可乘的（IV，第27页，命题2）；此外我们有

\[
\prod_ {i \in I} (1 + X _ {i} T) = 1 + \sum_ {k \geqslant 1} s _ {k} T ^ {k},\tag{14}
\]

其中形式幂级数 \(s_k \in \mathbf{A}[[\mathbf{X}]]\) 由下式定义

\[
s _ {k} = \sum_ {\mathrm{H} \in \mathfrak {P} _ {k}} \left(\prod_ {i \in \mathrm{H}} \mathrm{X} _ {i}\right) \quad (k \geqslant 1)\tag{15}
\]

（我们用 \(\mathfrak{P}_k\) 表示 I 的所有 \(k\) 元子集组成的集合）。特别地，我们有 \(s_1 = \sum_{i \in I} X_i\) 。当 I 是含 \(n\) 个元素的有限集时，我们有 \(s_k = 0\) 对 \(k > n\) 成立；更精确地说，当 \(\mathbf{I} = \{1,\dots,n\}\) 时，形式幂级数 \(s_k\) 正是 \(k\) 次的初等对称多项式（在 \(\mathbf{X}_1,\ldots,\mathbf{X}_n\) 中）。

设 \(\mathbf{S} = (\mathbf{S}_k)_{k\geqslant 1}\) 为一个未定元序列。由于形式幂级数 \(s_k\) 的阶 \(\geqslant k\) ，且属于 \(\mathbf{A}[[\mathbf{X}]]^{\mathrm{sym}}\) ，IV第28页命题4的条件 a) 和 b) 在 \(\mathrm{E} = \mathbf{A}[[\mathbf{X}]]^{\mathrm{sym}}\) 下得到满足；因此存在唯一的连续 A-代数同态

\[
\varphi_ {I}: \mathrm{A} [ [ \mathrm{S} ] ] \rightarrow \mathrm{A} [ [ \mathrm{X} ] ] ^ {\text { sym }}
\]

使得 \(\varphi_{I}(S_{k}) = s_{k}\) 对每个整数 \(k\geqslant 1\) 成立。

定理2. --- a) 若 I 是含 n 个元素的有限集，则 \(\varphi_{I}\) 诱导出从 A{[}{[}S \(_{1}\) , \ldots，S \(_{n}\) {]}{]} 到 A{[}{[}X{]}{]} \(^{\text{sym}}\) 上的双连续同构。

\begin{enumerate}
\def\labelenumi{\alph{enumi})}
\setcounter{enumi}{1}
\tightlist
\item
  如果 I 是无限集， \(\varphi_{\mathrm{I}}\) 是从 A{[}{[}S{]}{]} 到 A{[}{[}X{]}{]} \(^{\mathrm{sym}}\) 上的双连续同构。
\end{enumerate}

在情形 a) 中，我们令 \(\mathbf{B} = \mathbf{A}[[\mathbf{S}_1, \ldots, \mathbf{S}_n]]\) ，并赋予该代数由 \(\mathbf{A}[[\mathbf{S}]]\) 的拓扑诱导的拓扑；我们还用 \(\psi_{\mathrm{I}}\) 表示 \(\varphi_{\mathrm{I}}\) 在 \(\mathbf{B}\) 上的限制。在情形 b) 中，我们令 \(\mathbf{B} = \mathbf{A}[[\mathbf{S}]]\) ，且 \(\psi_{\mathrm{I}} = \varphi_{\mathrm{I}}\) 。我们将给多项式代数 \(\mathbf{A}[\mathbf{S}]\) 配备 \(\mathbf{N}\) 型分次，使得 \(\mathbf{S}_k\) 的权为 \(k\) （对每个整数 \(k \geqslant 1\) ）。

引理 1. --- 设 J 是 I 的有限子集，r 是使 Card \(J \geqslant r\) 成立的整数，且 f 是关于 \(X_{i}\) ( \(i \in I\) ) 的 r 次齐次对称形式幂级数。设 \(\bar{f}\) 为对 I - J 中每个 i 把 0 代入 \(X_{i}\) 所得的形式幂级数。如果 \(\bar{f} = 0\) ，则我们有 f = 0。

我们令 \(f = \sum_{|\alpha| = r} a_{\alpha} X^{\alpha}\) （其中 \(|\alpha|\) 是长度 \(\sum_{i \in I} \alpha_i\) ，即多重指标 \(\alpha = (\alpha_i)_{i \in I}\) 的长度）。如果 \(\alpha\) 是长度为 \(r\) 的多重指标，且 \(J'\) 是 \(\alpha\) 的支集（即由这样的 \(i \in I\) 组成的集合： \(\alpha_i \neq 0\) ），则 Card \(J' \leqslant r\) ，因此存在一个置换 \(\sigma \in \mathfrak{S}_I\) 使得 \(\sigma(J') \subset J\) 。令 \(\beta_i = \alpha_{\sigma^{-1}(i)}\) （对 \(i \in I\) ），那么单项式 \(X^\beta = \prod_{i \in I} X_{\sigma(i)}^{\alpha_i}\) 仅依赖于不定元 \(X_j (j \in J)\) ，从而 \(a_\beta = 0\) （由假设 \(\bar{f}=0\) ）。由于 f 是对称的，我们有 \(a_{\alpha}=a_{\beta}\) ，且由于 \(\alpha\) 是任意的，我们得出结论 f=0。

引理2. --- 设 \(f\) 是次数为 \(r\) 的、关于 \(X_i (i \in I)\) 的对称形式幂级数。存在唯一的多项式 \(P \in B \cap A[S]\) ，它是权为 \(r\) 的等权多项式，使得 \(f = \psi_I(P)\) 。

当I为有限时，该情形由定理1（IV，第62页）得出。

假设 I 是无限集，并选取 I 的一个含 r 个元素的有限子集 J。我们沿用引理 1 的记号；我们注意到， \(S_{n}\) ( \(n \geqslant 1\) ) 中每个权为 r 的等权多项式仅依赖于 \(S_{1}, \ldots, S_{r}\) ，且 \(\bar{s}_{1}, \ldots, \bar{s}_{r}\) 是 r 个不定元 \(X_{j}\) ( \(j \in J\) ) 中的初等对称多项式。若 P 是 \(S_{n}\) 中权为 r 的等权多项式，且 \(h = f - \psi_{\mathrm{I}}(\mathrm{P})\) ，则 \(\bar{h} = \bar{f} - \mathrm{P}(\bar{s}_{1}, \ldots, \bar{s}_{r})\) ，且引理 1 表明关系 \(f = \psi_{\mathrm{I}}(\mathrm{P})\) 等价于 \(\bar{f} = \mathrm{P}(\bar{s}_{1}, \ldots, \bar{s}_{r})\) 。由定理 1（IV，第 62 页），存在唯一的多项式 \(P \in A[S]\) ，等权且权为 r ，使得 \(\bar{f} = \mathrm{P}(\bar{s}_{1}, \ldots, \bar{s}_{r})\) ，由此即得结论。

引理 3. --- 对每个整数 \(m \geqslant 0\) ，令 \(\mathfrak{c}_m\) 为代数 \(\mathrm{A}[[\mathrm{X}]]^{\mathrm{sym}}\) 的理想，它由所有阶 \(\geqslant m\) 的对称形式幂级数组成。序列 \((\mathfrak{c}_m)_{m \geqslant 0}\) 是 \(\mathrm{A}[[\mathrm{X}]]^{\mathrm{sym}}\) 中 0 的邻域基。

当 I 为有限集时，引理显然成立，因此我们假设 I 为无限集。对 I 的每个由 m 个元素组成的有限子集 J，记 \(\tilde{J}\) 为 \(\mathbf{N}^{(1)}\) 中长度 \textless{} m 且支集含于 J 的元素组成的集合。进一步记 \(a_{j}^{\prime}\) 为不含任何形如 \(aX^{\alpha}\) （其中齐次 \(\alpha \in \tilde{J}\) ）的项的形式幂级数组成的集合。由于 \(\tilde{J}\) 是 \(\mathbf{N}^{(1)}\) 的有限子集，且 \(\mathbf{N}^{(1)}\) 的每个有限子集都包含于形如 \(\tilde{J}\) 的集合，族 \((\mathfrak{a}_{\mathrm{j}}^{\prime})\) 是 A{[}{[}X{]}{]} 中 0 的一个邻域基（IV，第26页）。现在，引理1蕴含关系式 \(a_{j}^{\prime} \cap A[[X]]^{\text{sym}} = c_{m}\) 对每个含 m 个元素的子集 J 成立，这便证明了引理3。

由于在 \(S_{k}\) 中给定权的单项式只有有限多个，每个形式幂级数 \(f \in B\) 都可以唯一地写成 \(f = \sum_{r \geqslant 0} P_{r}\) ，其中 \(P_{r}\) 是 \(B \cap A[S]\) 中权为 r 的等权多项式。对每个整数 \(m \geqslant 0\) ，令 \(\mathfrak{b}_m\) 为 B 中由所有上述类型的、满足 \(P_r = 0\) （对 \(0 \leqslant r < m\) ）的形式幂级数组成的理想。序列 \((\mathfrak{b}_m)_{m \geqslant 0}\) 是 B 中 0 的邻域基。

沿用上述记号， \(\psi_{\mathrm{I}}(\mathbf{P}_{r})\) 是关于 \(\mathbf{X}_i\) 的对称形式幂级数，且是 \(r\) 次齐次的，因而它是 \(r\) 次齐次分量（属于 \(\psi_{\mathrm{I}}(f)\) ）。引理2表明 \(\psi_{\mathrm{I}}\) 是 B 到 \(\mathbf{A}[[\mathbf{X}]]^{\mathrm{sym}}\) 上的代数同态，它把 \(\mathfrak{b}_m\) 映到 \(\mathfrak{c}_m\) （对每个整数 \(m \geqslant 0\) ）；于是引理3表明 \(\psi_{\mathrm{I}}\) 是双连续的，这就完成了定理2的证明。

